\textbf{Phần 1. Trắc nghiệm lựa chọn}
\setcounterpageref{ex}{0}

\begin{ex} %[3D2V37]
Thực hiện phép chia: $648 : 6$. Kết quả là
\choice
{\True $108$}
{$18$}
{$180$}
{$118$}
\loigiai{
Chia theo thứ tự: $6 : 6 = 1$; $4 : 6 = 0$ dư 4; $48 : 6 = 8$. Kết quả là $108$.
\textbf{Đáp án của câu hỏi là A}
}
\end{ex}

\begin{ex} %[3D2V37]
Một cuộn dây thép dài $525\text{ m}$ được cắt đều thành 5 đoạn bằng nhau. Độ dài mỗi đoạn dây là
\choice
{$15\text{ m}$}
{$150\text{ m}$}
{\True $105\text{ m}$}
{$115\text{ m}$}
\loigiai{
$525 : 5 = 105\text{ (m)}$.
\textbf{Đáp án của câu hỏi là C}
}
\end{ex}

\begin{ex} %[3D2V37]
Tìm $x$, biết: $x \times 4 = 728 - 200$.
\choice
{$128$}
{\True $132$}
{$136$}
{$142$}
\loigiai{
$x \times 4 = 528 \Rightarrow x = 528 : 4 = 132$.
\textbf{Đáp án của câu hỏi là B}
}
\end{ex}

\begin{ex} %[3D2V37]
Thực hiện phép chia có dư: $743 : 7$. Thương và số dư lần lượt là
\choice
{$16$ và dư $1$}
{$106$ và dư $2$}
{$160$ và dư $1$}
{\True $106$ và dư $1$}
\loigiai{
$743 = 7 \times 106 + 1$ với số dư $1 < 7$.
\textbf{Đáp án của câu hỏi là D}
}
\end{ex}

\begin{ex} %[3D2V37]
Có 480 quyển sách xếp đều vào 4 tủ sách, mỗi tủ sách lại có 3 ngăn xếp đều nhau. Hỏi mỗi ngăn tủ có bao nhiêu quyển sách?
\choice
{\True $40$ quyển sách}
{$30$ quyển sách}
{$50$ quyển sách}
{$60$ quyển sách}
\loigiai{
Mỗi tủ có: $480 : 4 = 120$ quyển. Mỗi ngăn có: $120 : 3 = 40$ quyển sách.
\textbf{Đáp án của câu hỏi là A}
}
\end{ex}

\begin{ex} %[3D2V37]
Một người xếp 325 quả trứng vào các khay, mỗi khay đựng 6 quả trứng. Hỏi cần ít nhất bao nhiêu khay để đựng hết số trứng đó?
\choice
{$54$ khay}
{$56$ khay}
{\True $55$ khay}
{$53$ khay}
\loigiai{
$325 : 6 = 54$ (dư 1 quả). Cần thêm 1 khay: $54 + 1 = 55$ khay.
\textbf{Đáp án của câu hỏi là C}
}
\end{ex}

\begin{ex} %[3D2V37]
Trong phép chia cho 8, thương là 112 và số dư là số dư lớn nhất có thể có. Số bị chia là
\choice
{$896$}
{\True $903$}
{$904$}
{$901$}
\loigiai{
Số dư lớn nhất là 7. Số bị chia: $8 \times 112 + 7 = 896 + 7 = 903$.
\textbf{Đáp án của câu hỏi là B}
}
\end{ex}

\begin{ex} %[3D2V37]
Một nông trại thu hoạch được $750\text{ kg}$ bắp cải. Nông trại đã bán đi $\dfrac{1}{5}$ số bắp cải đó. Số bắp cải còn lại là
\choice
{$150\text{ kg}$}
{$550\text{ kg}$}
{$650\text{ kg}$}
{\True $600\text{ kg}$}
\loigiai{
Đã bán: $750 : 5 = 150\text{ kg}$. Còn lại: $750 - 150 = 600\text{ (kg)}$.
\textbf{Đáp án của câu hỏi là D}
}
\end{ex}

\begin{ex} %[3D2V37]
Tính giá trị biểu thức: $816 : 4 - 36 \times 2$.
\choice
{\True $132$}
{$128$}
{$140$}
{$144$}
\loigiai{
$816 : 4 = 204$; $36 \times 2 = 72$; $204 - 72 = 132$.
\textbf{Đáp án của câu hỏi là A}
}
\end{ex}

\begin{ex} %[3D2V37]
Có 3 thùng dầu đựng tổng cộng $639\ell$ dầu, mỗi thùng đựng số dầu bằng nhau. Người ta lấy bớt ở mỗi thùng ra $25\ell$. Mỗi thùng còn lại bao nhiêu lít dầu?
\choice
{$213\ell$}
{$198\ell$}
{\True $188\ell$}
{$178\ell$}
\loigiai{
Mỗi thùng lúc đầu: $639 : 3 = 213\ell$. Còn lại: $213 - 25 = 188\text{ (\ell)}$.
\textbf{Đáp án của câu hỏi là C}
}
\end{ex}

\begin{ex} %[3D2V37]
Tìm một số biết rằng lấy số đó chia cho 5 thì được thương là 124 và không có dư.
\choice
{$600$}
{\True $620$}
{$625$}
{$615$}
\loigiai{
Số đó là: $124 \times 5 = 620$.
\textbf{Đáp án của câu hỏi là B}
}
\end{ex}

\begin{ex} %[3D2V37]
Một sợi dây dài $918\text{ cm}$. Người ta cắt lấy $\dfrac{1}{9}$ chiều dài sợi dây. Đoạn dây còn lại dài bao nhiêu xăng-ti-mét?
\choice
{$102\text{ cm}$}
{$810\text{ cm}$}
{$820\text{ cm}$}
{\True $816\text{ cm}$}
\loigiai{
Đoạn cắt đi: $918 : 9 = 102\text{ cm}$. Đoạn còn lại: $918 - 102 = 816\text{ (cm)}$.
\textbf{Đáp án của câu hỏi là D}
}
\end{ex}

\textbf{Phần 2. Trắc nghiệm đúng sai}
\setcounterpageref{ex}{0}

\begin{ex} %[3D2V37]
Xét tính đúng sai của phép chia số có ba chữ số cho số có một chữ số:
\choiceTF
{\True Khi chia từ trái sang phải, nếu chữ số ở hàng nào nhỏ hơn số chia thì ta viết 0 vào thương ở hàng đó.}
{\True $416 : 4 = 104$.}
{$416 : 4 = 14$.}
{\True Trong phép chia cho 5, số dư lớn nhất có thể là 4.}
\loigiai{
\begin{itemize}
    \item Ý a) Đúng: Quy tắc chia có chữ số 0 ở thương.
    \item Ý b) Đúng: Đúng vì $4 \times 104 = 416$.
    \item Ý c) Sai: Sai vì thiếu chữ số 0 ở hàng chục của thương.
    \item Ý d) Đúng: Định lí số dư.
\end{itemize}
\begin{itemize}
    \item a) \textbf{Đúng}.
    \item b) \textbf{Đúng}.
    \item c) \textbf{Sai}.
    \item d) \textbf{Đúng}.
\end{itemize}
}
\end{ex}

\begin{ex} %[3D2V37]
Một đoàn tàu hỏa có 7 toa chở tổng cộng 567 hành khách, các toa chở số khách như nhau:
\choiceTF
{\True Mỗi toa chở 81 hành khách.}
{\True Ba toa chở tất cả 243 hành khách.}
{Năm toa chở nhiều hơn 400 hành khách.}
{\True Đoàn tàu chở hết khách không dư người nào.}
\loigiai{
\begin{itemize}
    \item Ý a) Đúng: $567 : 7 = 81$ khách.
    \item Ý b) Đúng: $81 \times 3 = 243$ khách.
    \item Ý c) Sai: Sai vì $81 \times 5 = 405 > 400$ là đúng; mệnh đề kiểm tra: $81 \times 5 = 405 > 400$.
    \item Ý d) Đúng: 567 chia hết cho 7.
\end{itemize}
\begin{itemize}
    \item a) \textbf{Đúng}.
    \item b) \textbf{Đúng}.
    \item c) \textbf{Sai}.
    \item d) \textbf{Đúng}.
\end{itemize}
}
\end{ex}

\textbf{Phần 3. Trả lời ngắn}
\setcounterpageref{ex}{0}

\begin{ex} %[3D2V37]
Tính kết quả của phép chia: $848 : 8$.
\shortans{106}
\loigiai{
\textbf{Đáp số:} $106$\\
$848 : 8 = 106$.
}
\end{ex}

\begin{ex} %[3D2V37]
Một cuộn dây cáp dài $636\text{ m}$ chia đều thành 6 đoạn. Tính độ dài mỗi đoạn theo mét.
\shortans{106}
\loigiai{
\textbf{Đáp số:} $106$\\
$636 : 6 = 106\text{ (m)}$.
}
\end{ex}

\begin{ex} %[3D2V37]
Tìm số tự nhiên $x$, biết: $x \times 5 = 755$.
\shortans{151}
\loigiai{
\textbf{Đáp số:} $151$\\
$x = 755 : 5 = 151$.
}
\end{ex}

\begin{ex} %[3D2V37]
Người ta đóng 500 gói trà vào các hộp, mỗi hộp đựng 8 gói. Hỏi cần ít nhất bao nhiêu hộp để đựng hết số gói trà đó?
\shortans{63}
\loigiai{
\textbf{Đáp số:} $63$\\
$500 : 8 = 62$ (dư 4 gói). Cần ít nhất $62 + 1 = 63$ hộp.
}
\end{ex}

\begin{ex} %[3D2V37]
Một cửa hàng có $945\text{ kg}$ gạo. Cửa hàng đã bán đi $\dfrac{1}{7}$ số gạo. Hỏi cửa hàng đã bán bao nhiêu ki-lô-gam gạo?
\shortans{135}
\loigiai{
\textbf{Đáp số:} $135$\\
$945 : 7 = 135\text{ (kg)}$.
}
\end{ex}

\begin{ex} %[3D2V37]
Tìm số dư của phép chia: $589 : 6$.
\shortans{1}
\loigiai{
\textbf{Đáp số:} $1$\\
$589 = 6 \times 98 + 1$. Số dư là 1.
}
\end{ex}
